% Generated draft: requires mathematical and source review.
\documentclass[11pt]{article}
\usepackage[margin=25mm]{geometry}
\usepackage{fontspec}
\setmainfont{Latin Modern Roman}
\usepackage{amsmath,amssymb,mathtools,unicode-math}
\setmathfont{Latin Modern Math}
\usepackage{enumitem,needspace,xurl,hyperref}
\setlength{\parindent}{0pt}
\setlength{\parskip}{5pt}
\setlength{\emergencystretch}{3em}
\newcommand{\N}{\mathbb{N}}
\newcommand{\Z}{\mathbb{Z}}
\newcommand{\Q}{\mathbb{Q}}
\newcommand{\R}{\mathbb{R}}
\newcommand{\C}{\mathbb{C}}
\newcommand{\F}{\mathbb{F}}
\newcommand{\A}{\mathbb{A}}
\newcommand{\PP}{\mathbb{P}}
\newcommand{\E}{\mathbb{E}}
\newcommand{\eps}{\varepsilon}
\newcommand{\dd}{\,\mathrm{d}}
\begin{document}
\section*{UnsolvedMath ID: 30004290}
\textbf{Source problem code:} catalogue entry
\subsection{Definitions and mathematical statement}
For a positive integer \(n\), a divisor means a positive integer \(d\) satisfying \(d\mid n\), and \(\log\) denotes the natural logarithm. The supplied record asks for the typical asymptotic size of
\[
\Delta(n)=\max_i\#\{d\in\N:d\mid n,\ e^i\leq d\leq e^{i+1}\}.
\]
Both endpoints are included, as in the source record. The range of \(i\) is not specified. To expose this unresolved parameter, define, for \(I\in\{\Z,\R\}\),
\[
\Delta_I(n)=\max_{i\in I}\#\{d\in\N:d\mid n,\ e^i\leq d\leq e^{i+1}\}.
\]
Here \(I=\Z\) uses fixed intervals, whereas \(I=\R\) allows the interval to slide. These are distinct definitions; the supplied record does not determine which is intended.

A set \(A\subseteq\N\) has natural density one if
\[
\lim_{x\to\infty}\frac{\#\{n\in A:n\leq x\}}{\lfloor x\rfloor}=1.
\]
This gives a conventional interpretation of ``almost all integers.'' The record does not explicitly specify the probability or density convention underlying ``typical.''

\textbf{Question, subject to these qualifications.} Determine the asymptotic size of the source's \(\Delta(n)\) for a typical positive integer \(n\). The supplied curation additionally reports an almost-all lower bound
\[
\Delta(n)\gg(\log\log n)^{\eta+o(1)},
\qquad \eta\approx0.353,
\]
and asks whether its exponent is optimal. No exact definition of \(\eta\), range of the maximizing index, or precise quantification of the error term is supplied.

Under a natural-density interpretation, one possible precise reading of the reported lower bound is: for every \(\varepsilon>0\), there exists \(c_\varepsilon>0\) such that
\[
\lim_{x\to\infty}\frac{1}{\lfloor x\rfloor}
\#\left\{n\in\N:3\leq n\leq x,\ 
\Delta_I(n)\geq c_\varepsilon(\log\log n)^{\eta-\varepsilon}\right\}=1.
\]
This is an interpretation of the supplied claim, not an independently verified theorem. Under that interpretation, optimality could mean that no larger constant exponent satisfies the same almost-all lower-bound property. The stronger assertion that, for every \(\varepsilon>0\),
\[
\lim_{x\to\infty}\frac{1}{\lfloor x\rfloor}
\#\left\{n\in\N:3\leq n\leq x,\ 
(\log\log n)^{\eta-\varepsilon}\leq\Delta_I(n)
\leq(\log\log n)^{\eta+\varepsilon}\right\}=1
\]
would specify a typical exponent. The record does not distinguish these targets. The intended index domain, exact exponent, and meaning of optimality require clarification before a unique fully quantified formulation can be fixed.
\subsection{Short English statement}
The problem asks how many positive divisors of a typical integer can lie in a multiplicative interval of ratio \(e\), and whether the supplied almost-all lower-bound exponent \(\eta\approx0.353\) is optimal. The interval endpoints are inclusive. The maximizing index domain, exact value of \(\eta\), and intended meaning of optimality remain unresolved.
\subsection{Sources}
\begin{itemize}
\item Supplied local UnsolvedMath registry record, permanent ID \(30004290\), title ``Typical Size of the Erdős--Hooley Function,'' exact source problem code \texttt{catalogue entry}, in the Oberwolfach Reports Open Problems collection. The record lists proposed year 2019 and supplies no individual proposer.
\item Supplied report citation: \emph{Analytic Number Theory} (2020), Oberwolfach Reports, pp.~3141--3205, DOI \texttt{10.4171/owr/2019/50}. Original report link: \url{https://doi.org/10.4171/owr/2019/50}. The supplied background describes Oberwolfach Reports as documenting mathematical workshops and their open-problem discussions.
\item Supplied literature source: Ford--Green--Koukoulopoulos, ``equal sums and divisor concentration'' (descriptive label supplied by the registry): \url{https://arxiv.org/abs/1908.00378}.
\item Supplied literature source: Koukoulopoulos--Tao, ``upper mean-value bound'' (descriptive label supplied by the registry): \url{https://doi.org/10.1112/plms.12572}.
\item Literature assessment dated exactly \texttt{2026-08-22}: the supplied curation labels the problem \texttt{partially\_solved} and reports substantial progress on the general distribution question, sharp-looking lower bounds and extensive structural analysis of the Erdős--Hooley function, while stating that no complete description of the divisors of a typical integer is known. This is supplied dated curation derived from the OWR record, not an independently verified assessment; its time-sensitive status and cited claims require independent verification.
\item The registry labels the statement \texttt{corrected\_verified} and reports that the official Oberwolfach report was checked, the question was isolated from surrounding exposition, and notation was repaired. These are supplied provenance claims; no source was independently consulted for this draft.
\end{itemize}
\end{document}
